\section{Serving}
\label{sec:serving}

\subsection{Setup}

We measured \expone{} on one NVIDIA B300 GPU of an eight-GPU node: BF16
weights with the LoRA merged, transformers 5.17 with SDPA attention and flash-linear-attention
kernels for the Gated DeltaNet layers, Torch 2.12.1. The \texttt{causal\_conv1d} kernel was
not installed, and we used no \texttt{torch.compile}, CUDA graphs or quantization. Each
timing is measured warm: three untimed calls, then 20 timed repetitions (10 for the 8K and
16K inputs and for W4). HTTP timings use an in-process ASGI test client, so they include
request parsing, compilation, the model and reply validation, but no network. Each profile
of W1 to W3 is a single fixed request repeated; W4 uses 32 distinct requests. Our
specification asks for at least 200 distinct requests per profile, which this measurement
does not meet, and a 95th percentile from 20 samples is close to the second-largest value.
The workload question sets are listed in Appendix~\ref{app:prompt}.

\begin{table}[t]
\centering
\small
\caption{Single-B300 latency of \expone{} (BF16, merged LoRA). p50 / p95 in milliseconds.
``Model only'' excludes HTTP handling. W4 runs independent 1,024-token requests as one static
batch in one process.}
\label{tab:serving}
\setlength{\tabcolsep}{4pt}
\begin{adjustbox}{max width=\linewidth}
\begin{tabular}{@{}llrr@{}}
\toprule
Profile & Input & HTTP p50 / p95 & Model-only p50 \\
\midrule
W1 & 512 tokens, 8 candidates, 1 question & \textbf{97 / 98} & 96.5 \\
W2, full & 3.2k-token state, 8 questions (25.5k tokens) & 1,633 / 1,641 & 1,600 \\
W2, shared prefix & same request, 2,969-token prefix once & \textbf{379 / 380} & 343 \\
W3, 8K & 8,192 tokens, 1 question & 559 / 560 & 542 \\
W3, 16K & 16,320 tokens, 1 question & 1,166 / 1,169 & 1,139 \\
W3, K77 & 77 title candidates (2.1k tokens) & 147 / 148 & 141 \\
W3, K255 & 255 title candidates (6.6k tokens) & 447 / 449 & 433 \\
\midrule
 & & Batch p50 (ms) & Tokens/s \\
\cmidrule(l){3-4}
W4, batch 1 & 1,024 tokens & 93 & 11.0k \\
W4, batch 8 & 8,191 tokens & 515 & 15.9k \\
W4, batch 16 & 16,383 tokens & 1,019 & 16.1k \\
W4, batch 32 & 32,766 tokens & 2,054 & 16.0k \\
\bottomrule
\end{tabular}
\end{adjustbox}
\end{table}

\begin{figure}[t]
\centering
\figfile[width=0.9\linewidth]{figures/latency.pdf}
\caption{Left: HTTP p50 (bar) and p95 (dot) for each profile; the dashed lines mark the W1
latency targets of 150\,ms (p50) and 300\,ms (p95), which apply to W1 only. Right: W4
throughput against static batch size, labelled with the batch p50. One B300 GPU, BF16, merged
LoRA; warm, 20 timed repetitions (10 for W3 8K/16K and W4); in-process HTTP client (no
network); one fixed request per profile for W1 to W3.}
\label{fig:latency}
\end{figure}

\subsection{Results}

\paragraph{Short requests.}
W1 (a 512-token request with eight candidates and one question) takes 97\,ms at the median
and 98\,ms at the 95th percentile (Table~\ref{tab:serving}, Figure~\ref{fig:latency}), within
the specification's W1 targets of 150\,ms and 300\,ms. The formal gate G7 is defined for a
deployment that has passed the quality gates, which is not yet established. A 1,024-token
request (93\,ms) takes about as long as a 512-token request, so at this length fixed
per-kernel overhead dominates; CUDA graphs or compilation may reduce it.

\paragraph{Shared prefix.}
For W2 (a 3.2k-token state and eight questions), prefilling the 2,969-token common prefix
once and branching the hybrid cache (Section~\ref{sec:branching}) reduces the HTTP median from
1,633\,ms to 379\,ms, a factor of 4.3. Section~\ref{sec:determinism} reports its numerical
agreement with full sequences.

\paragraph{Long inputs and many candidates.}
A 16,320-token question takes 1.17\,s at the median; 255 title candidates (6.6k tokens) take
447\,ms. Peak memory was 56.5\,GB; Section~\ref{sec:engine} serves the same model on an 80\,GB
H100. The model is also larger than the 1 to 8B dense target of our specification.

\subsection{What was not measured}

W4 is a static batch in one process, not a server with a queue under concurrent load; we
measured neither queueing latency nor sustained concurrency. Timings exclude the network. We
did not measure unmerged serving, CUDA graphs, compilation or \texttt{causal\_conv1d} on this
path; Section~\ref{sec:engine} measures an 80\,GB H100 with vLLM, FP8 quantization and prefix
caching. The first run of these benchmarks
was lost with its host (Section~\ref{sec:ops}); values seen on screen during that run agreed
with the rerun reported here to within about 2\%.
